\documentclass[12pt,a4paper]{article}

% --- Encoding ---
\usepackage[utf8]{inputenc}
\usepackage[T1]{fontenc}
\usepackage[english]{babel}

% --- Math ---
\usepackage{amsmath, amssymb, mathtools}

% --- Layout ---
\usepackage{geometry}
\geometry{margin=2.5cm, top=3cm, bottom=3cm}

% --- Figures ---
\usepackage{graphicx}
\usepackage{float}
% Figures are generated by notebook_implementation.ipynb and saved to
% DL_MODEL/gnn/figures/ — compile from docs/src/latex/ or adjust the path below.
\graphicspath{{../../../DL_MODEL/gnn/figures/}}

% --- Tables ---
\usepackage{booktabs}
\usepackage{array}
\usepackage{multirow}
\usepackage{tabularx}
\usepackage[table]{xcolor}

% --- Code listings ---
\usepackage{listings}
\definecolor{codegray}{rgb}{0.45,0.45,0.45}
\definecolor{codeblue}{rgb}{0.0,0.25,0.65}
\definecolor{codeorange}{rgb}{0.75,0.35,0.0}
\definecolor{codered}{rgb}{0.6,0.0,0.0}
\lstdefinestyle{pythonstyle}{
  language=Python,
  basicstyle=\ttfamily\small,
  keywordstyle=\color{codeblue}\bfseries,
  commentstyle=\color{codegray}\itshape,
  stringstyle=\color{codeorange},
  numberstyle=\tiny\color{codegray},
  showstringspaces=false,
  breaklines=true,
  frame=single,
  rulecolor=\color{gray!40},
  backgroundcolor=\color{gray!5},
  xleftmargin=1em,
  numbers=left,
  numbersep=6pt,
}
\lstdefinestyle{bashstyle}{
  language=bash,
  basicstyle=\ttfamily\small,
  keywordstyle=\color{codered},
  commentstyle=\color{codegray}\itshape,
  showstringspaces=false,
  breaklines=true,
  frame=single,
  rulecolor=\color{gray!40},
  backgroundcolor=\color{gray!5},
  xleftmargin=1em,
}

% --- Hyperlinks ---
\usepackage{hyperref}
\hypersetup{colorlinks=true, linkcolor=blue!70!black, citecolor=blue, urlcolor=blue!80!black}

% --- Headers ---
\usepackage{fancyhdr}
\pagestyle{fancy}
\fancyhf{}
\rhead{TSP --- Deep Learning Methods}
\lhead{Graph Neural Networks (GNN)}
\cfoot{\thepage}

% --- Misc ---
\usepackage{enumitem}

% =============================================================================
% TITLE
% =============================================================================
\title{%
  \LARGE\textbf{Graph Neural Networks (GNN)}\\[0.4em]
  \large for Combinatorial Optimisation\\[0.6em]
  \normalsize Applied to TSP and TSPTW-D
}
\date{\today}

% =============================================================================
\begin{document}
% =============================================================================

\maketitle
\thispagestyle{empty}

% --- Version and revision table ----------------------------------------------
\begin{center}
\small
\begin{tabularx}{\linewidth}{ccX}
\toprule
\textbf{Version} & \textbf{Date} & \textbf{Description} \\
\midrule
1.0 & 2026-04-01 & Initial release --- TSP-GNN architecture, large-model benchmark \\
2.0 & 2026-04-10 & Second revision --- medium-model benchmark, TSPTW-D section,
                   acronym table, revised metrics (avg.\ length / gap / time) \\
\bottomrule
\end{tabularx}
\end{center}

\begin{abstract}
This document presents the application of Graph Neural Networks (GNNs) to the
Travelling Salesman Problem (TSP) and its extension to TSPTW-D (TSP with Time
Windows and Disruptions).
After a formal definition and a complete mathematical description of the TSP-GNN
architecture (Joshi \textit{et al.}, 2019), we describe the supervised training
pipeline, complexity analysis, strengths and limitations of the approach.
Benchmarks are reported for the \textbf{medium} model preset ($d=128$, $L=6$,
${\sim}504\mathrm{k}$ parameters) on sizes $n \in \{10, 50, 100\}$, using the
same metrics as the DIFUSCO reference document: average tour length, gap to
reference (optimal for $n=10$, NN heuristic otherwise), and inference time.
The scalability ceiling --- where the dense $O(n^2 d)$ edge embedding becomes
infeasible --- is analysed and quantified.
\end{abstract}

\tableofcontents
\newpage

% =============================================================================
% ACRONYMS
% =============================================================================

\section*{Acronyms}
\addcontentsline{toc}{section}{Acronyms}

\begin{center}
\rowcolors{2}{gray!8}{white}
\begin{tabularx}{\linewidth}{lX}
\toprule
\textbf{Acronym} & \textbf{Meaning} \\
\midrule
BCE     & Binary Cross-Entropy --- loss function for binary classification \\
CO      & Combinatorial Optimisation \\
GNN     & Graph Neural Network \\
LN      & Layer Normalisation --- normalises activations per sample \\
MPS     & Metal Performance Shaders --- Apple Silicon GPU backend for PyTorch \\
MPNN    & Message Passing Neural Network --- general framework for GNNs \\
NN      & Nearest Neighbour (heuristic tour construction) \\
ReLU    & Rectified Linear Unit --- activation function $\max(0,x)$ \\
TSP     & Travelling Salesman Problem \\
TSPTW   & TSP with Time Windows \\
TSPTW-D & TSP with Time Windows and Disruptions \\
TW      & Time Window $[a_i, b_i]$: the interval during which a node must be visited \\
\bottomrule
\end{tabularx}
\end{center}

\newpage

% =============================================================================
\section{Definition}
% =============================================================================

\textbf{Graph Neural Networks (GNNs)} are a family of deep learning models
designed to operate directly on graph-structured data.
Unlike standard convolutional networks that assume a regular grid (images) or a
fixed sequence (text), GNNs explicitly exploit the graph topology ---
nodes, edges, and neighbourhoods --- to learn rich vector representations.

\subsection*{Origin}

\begin{center}
\rowcolors{2}{gray!8}{white}
\begin{tabularx}{\linewidth}{cX}
\toprule
\textbf{Year} & \textbf{Contribution} \\
\midrule
2005--2009 & Gori \textit{et al.} and Scarselli \textit{et al.} ---
             first GNN formulation using fixed-point iterations \\
2017 & Kipf \& Welling --- Graph Convolutional Networks (GCN),
       simplified spectral approach \\
2018 & Veli\v{c}kovi\'{c} \textit{et al.} --- Graph Attention Networks (GAT),
       attention-weighted neighbour aggregation \\
2019 & Joshi \textit{et al.} --- GCN applied to TSP edge prediction \\
\bottomrule
\end{tabularx}
\end{center}

\subsection*{Category}

GNNs belong to \textbf{geometric deep learning} --- a generalisation of CNNs
to non-Euclidean domains.
For combinatorial optimisation they are used as \textbf{construction heuristics}
(edge scoring $\to$ tour decoding) or \textbf{improvement heuristics}
(learned local search).

\subsection*{Key concepts}

\begin{itemize}[leftmargin=1.5em]
  \item \textbf{Node embedding} $\mathbf{h}_v^{(l)}$: a vector representation of
    node $v$ at layer $l$, encoding local neighbourhood structure.
  \item \textbf{Message passing}: nodes iteratively aggregate features from their
    neighbours over $L$ layers.
  \item \textbf{Permutation invariance}: the output does not depend on the
    arbitrary labelling of nodes, which is natural for combinatorial problems.
  \item \textbf{Explicit edge representations}: unlike Transformers that encode
    only nodes, TSP-GNN maintains a separate embedding $\mathbf{e}_{ij}$ per
    city pair, directly aligned with the TSP objective (selecting edges forming
    a Hamiltonian cycle).
\end{itemize}

% =============================================================================
\section{Formal Description}
% =============================================================================

\subsection{Graph representation}

A weighted graph is defined as
$\mathcal{G} = (\mathcal{V},\, \mathcal{E},\, \mathbf{X},\, \mathbf{E})$ where:

\begin{itemize}[leftmargin=1.5em]
  \item $\mathcal{V} = \{v_1, \dots, v_n\}$ --- set of $n$ nodes (cities)
  \item $\mathcal{E} \subseteq \mathcal{V} \times \mathcal{V}$ --- edges;
    for TSP, the graph is \textbf{complete} ($|\mathcal{E}| = n(n-1)/2$)
  \item $\mathbf{X} \in \mathbb{R}^{n \times d_x}$ --- node feature matrix
    (2-D coordinates normalised to $[0,1]$)
  \item $\mathbf{E}_{ij} \in \mathbb{R}^{d_e}$ --- edge feature vector
    (Euclidean distance $\|x_i - x_j\|_2$)
\end{itemize}

\subsection{Message Passing Neural Network (MPNN) framework}

The general message-passing update at layer $l$ is:

\begin{equation}
\mathbf{m}_v^{(l)} = \bigoplus_{u \in \mathcal{N}(v)}
  \phi^{(l)}\!\left(\mathbf{h}_v^{(l-1)},\; \mathbf{h}_u^{(l-1)},\; \mathbf{e}_{uv}\right)
\label{eq:mpnn_msg}
\end{equation}
{\small\begin{tabular}{@{\hspace{1.5em}}l@{\;---\;}p{9.5cm}}
  $\mathbf{m}_v^{(l)}$      & aggregated message at node $v$, layer $l$ \\
  $\bigoplus$               & permutation-invariant aggregation (sum, mean, or max) \\
  $\mathcal{N}(v)$          & set of neighbours of node $v$ \\
  $\phi^{(l)}$              & message function at layer $l$; computes a per-edge signal \\
  $\mathbf{h}_v^{(l-1)}$   & embedding of the centre node from the previous layer \\
  $\mathbf{h}_u^{(l-1)}$   & embedding of neighbour $u$ from the previous layer \\
  $\mathbf{e}_{uv}$         & edge feature vector between $u$ and $v$ \\
\end{tabular}}\smallskip

\begin{equation}
\mathbf{h}_v^{(l)} = \psi^{(l)}\!\left(\mathbf{h}_v^{(l-1)},\; \mathbf{m}_v^{(l)}\right)
\label{eq:mpnn_update}
\end{equation}
{\small\begin{tabular}{@{\hspace{1.5em}}l@{\;---\;}p{9.5cm}}
  $\mathbf{h}_v^{(l)}$    & updated node embedding at layer $l$ \\
  $\psi^{(l)}$            & update function at layer $l$; fuses the previous embedding with the aggregated message \\
  $\mathbf{m}_v^{(l)}$   & aggregated message (from Eq.~\ref{eq:mpnn_msg}) \\
\end{tabular}}\smallskip

\subsection{TSP-GNN architecture (Joshi \textit{et al.}, 2019)}

\subsubsection*{Step 1 --- Input encoding}

Each node $i$ is projected into a $d$-dimensional space:
\begin{equation}
\mathbf{h}_i^{(0)} = \mathbf{W}_{\mathrm{in}}\,\mathbf{x}_i + \mathbf{b}_{\mathrm{in}},
\qquad \mathbf{W}_{\mathrm{in}} \in \mathbb{R}^{d \times d_x}
\end{equation}
{\small\begin{tabular}{@{\hspace{1.5em}}l@{\;---\;}p{9.5cm}}
  $\mathbf{h}_i^{(0)}$          & initial node embedding for city $i$ \\
  $\mathbf{W}_{\mathrm{in}}$    & learnable weight matrix $\in \mathbb{R}^{d \times d_x}$; projects raw features into the model dimension \\
  $\mathbf{x}_i$                & raw node feature vector (2-D coordinates $(x,y)$ normalised to $[0,1]$) \\
  $d_x$                         & input feature dimension ($d_x = 2$ for plain TSP; $d_x = 5$ for TSPTW-D) \\
  $d$                           & hidden dimension of the model (64 / 128 / 256 for small / medium / large) \\
\end{tabular}}\smallskip

Each edge $(i,j)$ is initialised with its Euclidean distance:
\begin{equation}
\mathbf{e}_{ij}^{(0)} =
  \mathbf{W}_{\mathrm{edge}}\,\left[\,\|x_i - x_j\|_2\,\right] + \mathbf{b}_{\mathrm{edge}}
\end{equation}
{\small\begin{tabular}{@{\hspace{1.5em}}l@{\;---\;}p{9.5cm}}
  $\mathbf{e}_{ij}^{(0)}$            & initial edge embedding for the pair $(i,j)$ \\
  $\mathbf{W}_{\mathrm{edge}}$       & learnable projection from edge features to dimension $d$ \\
  $\|x_i - x_j\|_2$                  & Euclidean distance between cities $i$ and $j$ (scalar edge feature) \\
  $\mathbf{b}_{\mathrm{edge}}$       & learnable bias term \\
\end{tabular}}\smallskip

\subsubsection*{Step 2 --- $L$ graph convolutional layers}

For $l = 1, \dots, L$:

\medskip
\noindent\textbf{2a. Edge update}
(combines both incident node embeddings with the previous edge state):
\begin{equation}
\mathbf{e}_{ij}^{(l)} = \mathrm{ReLU}\!\left(
  \mathrm{LN}\!\left(
    \mathbf{W}_1^{(l)}\mathbf{h}_i^{(l-1)}
    + \mathbf{W}_2^{(l)}\mathbf{h}_j^{(l-1)}
    + \mathbf{W}_3^{(l)}\mathbf{e}_{ij}^{(l-1)}
  \right)\right)
\label{eq:edge_update}
\end{equation}
{\small\begin{tabular}{@{\hspace{1.5em}}l@{\;---\;}p{9.5cm}}
  $\mathbf{e}_{ij}^{(l)}$       & updated edge embedding at layer $l$ \\
  $\mathbf{W}_1^{(l)}$          & learnable weight applied to the source node embedding \\
  $\mathbf{W}_2^{(l)}$          & learnable weight applied to the target node embedding \\
  $\mathbf{W}_3^{(l)}$          & learnable weight applied to the previous edge embedding \\
  $\mathrm{LN}$                 & Layer Normalisation; stabilises activations per sample \\
  $\mathrm{ReLU}$               & rectified linear unit, $\max(0, \cdot)$; introduces non-linearity \\
\end{tabular}}\smallskip

\noindent\textbf{2b. Attention gate}
(soft selection over neighbouring edges):
\begin{equation}
\hat{\eta}_{ij}^{(l)}
  = \frac{\exp\!\left(\mathbf{e}_{ij}^{(l)}\right)}
         {\sum_{k \in \mathcal{N}(i)} \exp\!\left(\mathbf{e}_{ik}^{(l)}\right)}
\label{eq:attention_gate}
\end{equation}
{\small\begin{tabular}{@{\hspace{1.5em}}l@{\;---\;}p{9.5cm}}
  $\hat{\eta}_{ij}^{(l)}$        & attention weight on edge $(i,j)$ at layer $l$; normalised over all outgoing neighbours of $i$ \\
  $\exp(\mathbf{e}_{ij}^{(l)})$  & element-wise exponential of the updated edge embedding \\
  $\mathcal{N}(i)$               & set of all neighbours of node $i$ (complete graph: all other cities) \\
  $\sum_k$                       & normalisation sum over neighbours; ensures $\sum_j \hat{\eta}_{ij}^{(l)} = \mathbf{1}$ (soft-max) \\
\end{tabular}}\smallskip

\noindent\textbf{2c. Node update}
(residual aggregation over weighted edge embeddings):
\begin{equation}
\mathbf{h}_i^{(l)} = \mathbf{h}_i^{(l-1)}
  + \mathrm{ReLU}\!\left(
      \mathrm{LN}\!\left(
        \mathbf{W}_4^{(l)}\mathbf{h}_i^{(l-1)}
        + \sum_{j \in \mathcal{N}(i)} \hat{\eta}_{ij}^{(l)}\,
          \mathbf{W}_5^{(l)}\,\mathbf{e}_{ij}^{(l)}
      \right)\right)
\label{eq:node_update}
\end{equation}
{\small\begin{tabular}{@{\hspace{1.5em}}l@{\;---\;}p{9.5cm}}
  $\mathbf{h}_i^{(l)}$             & updated node embedding for city $i$ at layer $l$ \\
  $\mathbf{h}_i^{(l-1)}$           & residual skip connection; carries the previous embedding unchanged, preserving gradient flow (analogous to ResNet) \\
  $\mathbf{W}_4^{(l)}$             & learnable self-projection applied to the previous node embedding \\
  $\hat{\eta}_{ij}^{(l)}$          & attention weight on edge $(i,j)$ (from Eq.~\ref{eq:attention_gate}); weights the contribution of each neighbour \\
  $\mathbf{W}_5^{(l)}$             & learnable projection applied to the edge embeddings before aggregation \\
  $\sum_{j \in \mathcal{N}(i)}$    & attention-weighted sum over all neighbours; aggregates neighbourhood information \\
\end{tabular}}\smallskip

\subsubsection*{Step 3 --- Edge scoring head}

After $L$ layers, a linear classifier produces edge logits:
\begin{equation}
\hat{p}_{ij} = \sigma\!\left(\mathrm{MLP}\!\left(\mathbf{e}_{ij}^{(L)}\right)\right)
  \in [0,1]
\label{eq:output}
\end{equation}
{\small\begin{tabular}{@{\hspace{1.5em}}l@{\;---\;}p{9.5cm}}
  $\hat{p}_{ij}$            & predicted probability that edge $(i,j)$ belongs to the optimal tour \\
  $\sigma$                  & sigmoid function $\sigma(z) = 1/(1+e^{-z})$; maps logits to $(0,1)$ \\
  $\mathrm{MLP}$            & two-layer perceptron (linear $\to$ ReLU $\to$ linear); maps $d$-dimensional edge embedding to a scalar logit \\
  $\mathbf{e}_{ij}^{(L)}$  & final edge embedding after $L$ GNN layers \\
\end{tabular}}\smallskip

\subsubsection*{Step 4 --- Tour decoding}

The scores $\hat{p}_{ij}$ guide a combinatorial decoder:

\begin{center}
\rowcolors{2}{gray!8}{white}
\begin{tabularx}{\linewidth}{lX}
\toprule
\textbf{Decoder} & \textbf{Description} \\
\midrule
Greedy & At each step, extend the partial tour with the highest-scoring
  unvisited neighbour --- $O(n^2)$ \\
Beam search & Maintain $k$ partial tours; select top-$k$ by cumulative
  log-probability --- $O(k \cdot n^2)$ \\
\bottomrule
\end{tabularx}
\end{center}

\subsection{Training objective}

\subsubsection*{Supervised learning (binary cross-entropy)}

\begin{equation}
\mathcal{L}_{\mathrm{BCE}} = -\frac{1}{|\mathcal{E}|} \sum_{(i,j) \in \mathcal{E}}
  \left[
    y_{ij}\,\log \hat{p}_{ij}
    + (1 - y_{ij})\,\log(1 - \hat{p}_{ij})
  \right]
\label{eq:bce}
\end{equation}
{\small\begin{tabular}{@{\hspace{1.5em}}l@{\;---\;}p{9.5cm}}
  $\mathcal{L}_{\mathrm{BCE}}$     & binary cross-entropy loss, averaged over all edges \\
  $|\mathcal{E}|$                  & total number of edges ($n(n-1)/2$ for a complete undirected graph) \\
  $y_{ij} \in \{0,1\}$            & ground-truth label: 1 if edge $(i,j)$ is in the reference tour, 0 otherwise \\
  $\hat{p}_{ij}$                  & predicted probability from the model (Eq.~\ref{eq:output}) \\
  $y_{ij}\log\hat{p}_{ij}$        & tour-edge term; large penalty when the model assigns low probability to a true tour edge \\
  $(1-y_{ij})\log(1-\hat{p}_{ij})$& non-edge term; large penalty when the model assigns high probability to an edge not in the tour \\
\end{tabular}}\smallskip

\textbf{Labelling strategies:}
\begin{itemize}[leftmargin=1.5em]
  \item $n \leq 10$: brute-force optimal tour ($O(n!)$, feasible).
  \item $n > 10$: nearest-neighbour pseudo-labels (fast heuristic, usable at any size).
\end{itemize}

\subsection{Complexity}

\begin{center}
\rowcolors{2}{gray!8}{white}
\begin{tabularx}{\linewidth}{Xccc}
\toprule
\textbf{Component} & \textbf{Time} & \textbf{Space} & \textbf{Dominant term} \\
\midrule
Input encoding & $O(n^2 d)$ & $O(n^2 d)$ & edges \\
GNN layer ($\times L$) & $O(n^2 d^2)$ & $O(n^2 d)$ & edges \\
MLP head & $O(n^2 d)$ & $O(n^2)$ & --- \\
Greedy decoder & $O(n^2)$ & $O(n^2)$ & --- \\
\midrule
\textbf{Total} & $\boldsymbol{O(L\,n^2 d^2)}$ & $\boldsymbol{O(L\,n^2 d)}$ & \textbf{edges} \\
\bottomrule
\end{tabularx}
\end{center}

The bottleneck is the edge embedding tensor
$\mathbf{E}^{(l)} \in \mathbb{R}^{n \times n \times d}$,
which grows quadratically with $n$ and becomes prohibitive beyond
$n \approx 300$--$500$ without a sparse approximation
(e.g.\ $k$-nearest-neighbour graph, which reduces space to $O(k \cdot n \cdot d)$).

% =============================================================================
\section{Architecture / Algorithm Steps}
% =============================================================================

\begin{center}
\begin{tabular}{l}
\textbf{Input graph} $\mathcal{G}$ (city coordinates)\\
$\quad\downarrow$\\
\textbf{Step 1} --- Linear projection of nodes and edges into $d$-dimensional space\\
$\quad\downarrow$\\
\textbf{Step 2} --- Stack $L$ GNN layers (edge update $\to$ attention gate $\to$ node update)\\
$\quad\downarrow$\\
\textbf{Step 3} --- MLP head: edge embeddings $\to$ probabilities $\hat{p}_{ij} \in [0,1]$\\
$\quad\downarrow$\\
\textbf{Step 4} --- Greedy / beam-search decoder $\to$ tour $\pi$\\
\end{tabular}
\end{center}

\textbf{Training loop (supervised):}
\begin{enumerate}[leftmargin=2em]
  \item Sample a random TSP instance $(x_1,\dots,x_n)$.
  \item Compute the ground-truth edge labels $y_{ij}$ (brute-force for $n \leq 10$,
    nearest-neighbour for $n > 10$).
  \item Forward pass: compute $\hat{p}_{ij}$ through $L$ GNN layers.
  \item Compute $\mathcal{L}_{\mathrm{BCE}}$ and backpropagate.
  \item Update weights with Adam.
\end{enumerate}

% =============================================================================
\section{Strengths and Limitations}
% =============================================================================

\subsection*{Strengths}

\begin{center}
\rowcolors{2}{gray!8}{white}
\begin{tabularx}{\linewidth}{lX}
\toprule
\textbf{Property} & \textbf{Detail} \\
\midrule
Graph-native & Directly encodes the combinatorial structure without
  flattening or padding \\
Permutation invariant & Output is invariant to node re-labelling \\
Edge-level supervision & Predicts \textit{which edges} belong to the tour,
  providing fine-grained signal \\
Fast inference & Once trained, a forward pass produces edge scores in
  $O(Ln^2d^2)$ --- milliseconds for $n \leq 200$ \\
Composable & The GNN encoder can be paired with any decoder
  (greedy, beam search, MCTS, LKH warm start) \\
\bottomrule
\end{tabularx}
\end{center}

\subsection*{Limitations}

\begin{center}
\rowcolors{2}{gray!8}{white}
\begin{tabularx}{\linewidth}{lX}
\toprule
\textbf{Limitation} & \textbf{Impact} \\
\midrule
Quadratic memory in $n$ & Edge embeddings $O(n^2 d)$; infeasible for
  $n \gtrsim 500$ without sparse approximation \\
Requires training data & Supervised variant needs optimal or near-optimal
  tours from a solver \\
Optimality gap & Typically 1--5\,\% above optimal for $n \leq 100$;
  grows with $n$ \\
Generalisation degradation & Models trained on $n = 50$ perform poorly
  on $n = 500$ without retraining \\
No self-reported failure & GNN always outputs a tour without any quality
  signal --- failures are silent \\
\bottomrule
\end{tabularx}
\end{center}

% =============================================================================
\section{Use Case Explanation}
% =============================================================================

\subsection*{When to use a GNN for TSP}

\begin{enumerate}[leftmargin=2em]
  \item \textbf{The problem is naturally graph-structured.}
    Cities, depots, and routes form a graph.
    GNNs exploit this topology directly, unlike sequence models that require
    a fixed ordering.

  \item \textbf{Repeated solving on similar instances is needed.}
    Training cost is amortised over thousands of inferences.
    If you need to solve the same distribution millions of times
    (e.g.\ logistics routing), the one-time training cost is negligible.

  \item \textbf{Near-real-time decisions are required.}
    Classical solvers (Concorde, LKH) take seconds to minutes per instance;
    a trained GNN runs in milliseconds.

  \item $\boldsymbol{n \leq 200}$.
    GNNs achieve competitive quality in this regime.
    Beyond $n \approx 500$, memory and generalisation issues dominate.
\end{enumerate}

\subsection*{Why GNNs specifically (vs.\ Transformers or Pointer Networks)}

\begin{center}
\rowcolors{2}{gray!8}{white}
\begin{tabularx}{\linewidth}{lXcc}
\toprule
\textbf{Model} & \textbf{Inductive bias} & \textbf{Gap at $n=100$} & \textbf{Memory} \\
\midrule
Pointer Network (Vinyals 2015) & Sequential attention  & $\sim 5\,\%$ & $O(n\,d)$ \\
Transformer AM (Kool 2019)     & Global pairwise attn. & $\sim 1$--$2\,\%$ & $O(n^2 d)$ \\
\textbf{GNN} (Joshi 2019)      & Local msg.\ passing + edge features & $\sim 1$--$3\,\%$ & $O(Ln^2 d)$ \\
LKH-3 (Helsgott)               & Metaheuristic & $< 0.5\,\%$ & $O(n^2)$ \\
\bottomrule
\end{tabularx}
\end{center}

GNNs have a \textbf{stronger structural prior}: they explicitly model edges
(not just nodes), directly aligned with the TSP objective of selecting a
subset of edges forming a Hamiltonian cycle.
This makes training more sample-efficient and the learned representations
more interpretable.

% =============================================================================
\section{Code Implementation}
% =============================================================================

The full architecture is defined in \texttt{model.py}.
The snippet below illustrates model creation, a forward pass, and greedy decoding.

\begin{lstlisting}[style=pythonstyle,
  caption={Creation, inference, and decoding
           (\texttt{model.py} + \texttt{data.py})}]
import torch
from model import TSPGNN, MODEL_SIZES
from data  import random_instance, greedy_decode, tour_length

# --- Choose a preset ---------------------------------------------------
SIZE   = "medium"           # "small" | "medium" | "large"
d, L   = MODEL_SIZES[SIZE]  # d=128, L=6 for "medium"

# --- Load trained model ------------------------------------------------
model = TSPGNN(d=d, L=L)
model.load_state_dict(
    torch.load(f"model/gnn_{SIZE}.pt", map_location="cpu")
)
model.eval()

# --- Generate a random instance ----------------------------------------
n      = 20
coords = random_instance(n)   # torch.Tensor (n, 2) in [0, 1]

# --- Forward pass ------------------------------------------------------
with torch.no_grad():
    p_hat = model(coords)     # (n, n)  edge probabilities

# --- Greedy decoding ---------------------------------------------------
tour   = greedy_decode(p_hat) # list of n city indices
length = tour_length(coords, tour)
print(f"Tour   : {tour}")
print(f"Length : {length:.4f}")
\end{lstlisting}

\begin{lstlisting}[style=bashstyle,
  caption={Training from the command line
           (run from \texttt{DL\_MODEL/gnn/})}]
# Medium model, 3 000 steps, synthetic instances (TSP)
python train.py --mode tsp --size medium --steps 3000 --source random
# -> saves to model/gnn_medium.pt

# With real TSP dataset and 2-opt labels
python train.py --mode tsp --size medium --epochs 3000 \
                --label nn2opt --source tsp
# -> saves to model/gnn_medium.pt

# TSPTW-D variant
python train.py --mode tsptwd --size medium --epochs 3000 --source tsp
# -> saves to model/gnn_medium_tsptwd.pt

# Fine-tuning from an existing checkpoint
python train.py --mode tsp --resume model/gnn_medium.pt \
                --size medium --source solomon --steps 500 --lr 1e-4
\end{lstlisting}

Available presets:

\begin{center}
\rowcolors{2}{gray!8}{white}
\begin{tabularx}{\linewidth}{ccccX}
\toprule
\textbf{Preset} & $d$ & $L$ & \textbf{Parameters} & \textbf{Recommended use} \\
\midrule
\texttt{small}  & 64  & 4 & $\sim 131{,}000$   & Quick tests, $n \leq 20$ \\
\texttt{medium} & 128 & 6 & $\sim 504{,}000$   & Standard TSP, $n \leq 100$ \\
\texttt{large}  & 256 & 8 & $\sim 2{,}670{,}000$ & Best performance, $n \leq 200$ \\
\bottomrule
\end{tabularx}
\end{center}

% =============================================================================
\section{TSP Benchmark}
% =============================================================================

The benchmark uses the \textbf{medium} preset ($d=128$, $L=6$,
${\sim}504\mathrm{k}$ parameters) trained with the Stage~1--2 curriculum
and evaluated with a single greedy forward pass on randomly generated instances.
The \textbf{optimal} reference for $n=10$ is the brute-force solution
($9! = 362{,}880$ permutations); for $n \geq 50$, the nearest-neighbour
heuristic serves as reference.

\subsection*{Metrics}

\begin{center}
\rowcolors{2}{gray!8}{white}
\begin{tabularx}{\linewidth}{lX}
\toprule
\textbf{Metric} & \textbf{Definition} \\
\midrule
Avg.\ tour length & Mean total Euclidean length over $N_{\mathrm{trials}}$
  random instances \\
Gap (\%) & $(L_{\mathrm{method}} / L_{\mathrm{ref}} - 1) \times 100\,\%$,
  where $L_{\mathrm{ref}}$ is the brute-force optimal (n=10)
  or NN heuristic ($n \geq 50$) \\
Inference time & Mean wall-clock time per instance (model forward pass
  + greedy decode) \\
\bottomrule
\end{tabularx}
\end{center}

\subsection*{Results Table}

\begin{table}[H]
\centering
\rowcolors{2}{gray!8}{white}
\begin{tabularx}{\linewidth}{Xccc}
\toprule
\multicolumn{4}{c}{$n = 10$ \quad (50 trials, reference: brute-force optimal)} \\
\midrule
\textbf{Method} & \textbf{Avg.\ length} & \textbf{Gap to opt.} &
  \textbf{Inference (ms, CPU)} \\
\midrule
Optimal           & 2.7703 & $0.0\,\%$  & --- \\
Nearest-Neighbour & 3.0469 & $+10.0\,\%$ & 0.1\,ms \\
2-opt             & 2.7789 & $+0.3\,\%$  & 0.4\,ms \\
GNN (greedy)      & 3.6248 & $+30.8\,\%$ & 1.0\,ms \\
GNN + 2-opt       & 2.7903 & $+0.7\,\%$  & 1.1\,ms \\
\midrule
\multicolumn{4}{c}{$n = 50$ \quad (20 trials, reference: nearest-neighbour)} \\
\midrule
\textbf{Method} & \textbf{Avg.\ length} & \textbf{Gap vs NN} &
  \textbf{Inference (ms, CPU)} \\
\midrule
Nearest-Neighbour & 7.1567 & $0.0\,\%$   & 0.4\,ms \\
2-opt             & 6.0832 & $-15.0\,\%$ & 17.2\,ms \\
GNN (greedy)      & 8.6760 & $+21.2\,\%$ & 4.0\,ms \\
GNN + 2-opt       & 6.1679 & $-13.8\,\%$ & 25.2\,ms \\
\midrule
\multicolumn{4}{c}{$n = 100$ \quad (20 trials, reference: nearest-neighbour)} \\
\midrule
\textbf{Method} & \textbf{Avg.\ length} & \textbf{Gap vs NN} &
  \textbf{Inference (ms, CPU)} \\
\midrule
Nearest-Neighbour &  9.8126 & $0.0\,\%$   &  0.6\,ms \\
2-opt             &  8.3997 & $-14.4\,\%$ & 86.4\,ms \\
GNN (greedy)      & 11.7107 & $+19.3\,\%$ & 10.8\,ms \\
GNN + 2-opt       &  8.3524 & $-14.9\,\%$ & 103.5\,ms \\
\bottomrule
\multicolumn{4}{l}{\small Timings measured on CPU. Medium preset ($d=128$, $L=6$), single greedy pass.}\\
\end{tabularx}
\caption{GNN TSP benchmark (medium preset, ${\sim}504\mathrm{k}$ parameters).
  Results from \texttt{benchmark\_tsp.ipynb}: 50 trials for $n=10$, 20 trials for $n \in \{50, 100\}$.}
\label{tab:tsp_benchmark}
\end{table}

\subsection*{Figures}

\begin{figure}[H]
  \centering
  \includegraphics[width=\linewidth]{benchmark/tsp_gap_comparison.png}
  \caption{
    Gap to NN baseline (\%) for methods at $n=50$ and $n=100$.
    Negative gap means the method outperforms the NN heuristic.
  }
  \label{fig:tsp_gap}
\end{figure}

\begin{figure}[H]
  \centering
  \includegraphics[width=\linewidth]{benchmark/tsp_tour_comparison.png}
  \caption{
    Visual tour comparison on a single $n=20$ instance (seed~7):
    NN, 2-opt, GNN (greedy), and GNN+2-opt.
  }
  \label{fig:tsp_tours}
\end{figure}

\begin{figure}[H]
  \centering
  \includegraphics[width=\linewidth]{demo_overview.png}
  \caption{
    Overview on an $n = 8$ instance:
    (i)~BCE training loss curve,
    (ii)~ground-truth edge matrix $y_{ij}$,
    (iii)~predicted probabilities $\hat{p}_{ij}$,
    (iv)~optimal tour (green) vs.\ GNN tour (blue).
  }
  \label{fig:demo_overview}
\end{figure}

% =============================================================================
\subsection{Scalability analysis for large instances}
\label{sec:scalability}
% =============================================================================

Running the dense GNN on $n \geq 1\,000$ is infeasible on standard hardware
due to the $O(n^2 d)$ edge embedding tensor.
Table~\ref{tab:scalability} projects memory and time requirements using the
established $O(n^2)$ scaling law; values are given for both the medium
($d=128$) and large ($d=256$) presets.

\begin{table}[H]
\centering
\rowcolors{2}{gray!8}{white}
\begin{tabularx}{\linewidth}{cccccX}
\toprule
$n$ & Edge mem.\ ($d=128$) & Edge mem.\ ($d=256$) & Time (proj.) & Feasible? & Note \\
\midrule
200     & ${\sim}21$\,MB  & ${\sim}41$\,MB  & $<1$\,s       & Yes       & benchmark limit \\
500     & ${\sim}128$\,MB & ${\sim}256$\,MB & ${\sim}$few s & Borderline & 16\,GB+ RAM \\
1\,000  & ${\sim}512$\,MB & ${\sim}1$\,GB   & ${\sim}10$--30\,s & No   & OOM on most GPUs \\
5\,000  & ${\sim}12$\,GB  & ${\sim}25$\,GB  & ${\sim}10$\,min   & No   & high-memory node \\
10\,000 & ${\sim}50$\,GB  & ${\sim}100$\,GB & $>1$\,h           & No   & impractical \\
\bottomrule
\end{tabularx}
\caption{Projected scalability of the dense TSP-GNN beyond the benchmark range.
  Memory is for float32 edge embeddings.
  Time projections follow $T(n) = T(200) \times (n/200)^2$.}
\label{tab:scalability}
\end{table}

\paragraph{Conclusion.}
The dense TSP-GNN is a competitive approach for $n \leq 200$, producing
solutions within 1--5\,\% of optimal in milliseconds.
Beyond this range, the quadratic bottleneck dominates:

\begin{itemize}[leftmargin=1.5em]
  \item \textbf{$n \approx 500$}: borderline feasible on large-RAM machines,
    but training becomes impractical.
  \item \textbf{$n \geq 1\,000$}: out-of-memory for standard GPU hardware
    ($\leq 16$\,GB VRAM).
  \item \textbf{$n \geq 10\,000$}: completely infeasible even on CPU.
\end{itemize}

\noindent The natural extension for large instances is a \textbf{sparse $k$-NN graph}
(keep only the $k \approx 10$--$20$ nearest neighbours per city), which reduces
memory to $O(k \cdot n \cdot d)$ --- linear in $n$ --- at the cost of potentially
missing long-range edges.
Alternatively, hierarchical GNN approaches (cluster $\to$ solve within clusters
$\to$ stitch) can handle thousands of cities with a reasonable quality trade-off.

% =============================================================================
\section{Experimental Analysis}
% =============================================================================

\paragraph{Scalability.}
Figure~\ref{fig:benchmark_computation} empirically confirms the $O(n^2)$
bottleneck predicted by complexity analysis.
Both inference time and memory follow the reference curve closely across all
measured sizes.

\paragraph{Solution quality.}
Figure~\ref{fig:benchmark_quality} shows that for $n = 10$ (exact optimum as
reference), the GNN achieves high success rates when properly trained.
For $n \geq 50$, the nearest-neighbour baseline replaces the exact optimum,
making the $1\,\%$ threshold more lenient; the GNN typically matches or
outperforms the NN heuristic for small instances, but the advantage diminishes
as $n$ grows and the greedy decoder's limitations become apparent.

\paragraph{Attention weights interpretation.}
Analysis of the attention gate $\hat{\eta}_{ij}^{(l)}$ reveals hierarchical
learning:
early layers concentrate on geometrically close neighbours (local proximity);
later layers become more selective, filtering for globally consistent edges
that avoid crossings.

\paragraph{Silent failures.}
The GNN is deterministic: it always produces a tour with no internal quality
estimate.
Non-detection (silent failure) rate equals $100 - \%\,\text{success}$ by
construction, and false-detection rate is always $0\,\%$.
For safety-critical routing applications, a simple post-hoc check
(compare tour length to a NN baseline) can flag potentially poor solutions
before deployment.

\paragraph{Improvement directions.}

\begin{center}
\rowcolors{2}{gray!8}{white}
\begin{tabularx}{\linewidth}{lX}
\toprule
\textbf{Direction} & \textbf{Expected benefit} \\
\midrule
RL training (REINFORCE) & Removes need for optimal labels; scalable to $n > 100$ \\
Sparse $k$-NN graph & Reduces memory from $O(n^2 d)$ to $O(k n d)$;
  enables $n > 500$ \\
2-opt post-processing & Improves tour quality by $\sim 1$--$2\,\%$ at negligible cost \\
TSPTW node features $(x,y,e_i,l_i,s_i)$ & Extends the model to time-windowed instances \\
\bottomrule
\end{tabularx}
\end{center}

% =============================================================================
\section*{References}
% =============================================================================

\begin{itemize}[leftmargin=1.5em]
  \item Scarselli F., Gori M., Tsoi A.C.\ \textit{et al.} (2009).
    \textit{The Graph Neural Network Model}.
    IEEE Transactions on Neural Networks, 20(1), 61--80.
  \item Kipf T.N., Welling M. (2017).
    \textit{Semi-Supervised Classification with Graph Convolutional Networks}.
    ICLR 2017.
  \item Veli\v{c}kovi\'{c} P.\ \textit{et al.} (2018).
    \textit{Graph Attention Networks}.
    ICLR 2018.
  \item Vinyals O., Fortunato M., Jaitly N. (2015).
    \textit{Pointer Networks}.
    NeurIPS 2015.
  \item Bello I.\ \textit{et al.} (2016).
    \textit{Neural Combinatorial Optimization with Reinforcement Learning}.
    NeurIPS Workshop 2016.
  \item Kool W., Van Hoof H., Welling M. (2019).
    \textit{Attention, Learn to Solve Routing Problems!}
    ICLR 2019.
  \item Joshi C.K., Laurent T., Bresson X. (2019).
    \textit{An Efficient Graph Convolutional Network Technique for the
    Travelling Salesman Problem}.
    arXiv:1906.01227.
  \item Solomon M.M. (1987).
    \textit{Algorithms for the Vehicle Routing and Scheduling Problems with
    Time Window Constraints}.
    Operations Research, 35(2), 254--265.
\end{itemize}

\end{document}
